\documentclass[12pt]{article}
\usepackage[utf8]{inputenc}
\usepackage{geometry}
\geometry{a4paper, margin=1in}
\usepackage{booktabs}
\usepackage{graphicx}
\usepackage{caption}
\usepackage{float}
\usepackage{pdflscape}

\title{Extension: Continuous Cognitive Sensitivity Analysis}
\author{Research Team}
\date{\today}

\begin{document}
\maketitle

\section{Introduction and Research Question}
This document serves as a reproducible extension to the primary analysis. The core research question addressed is: \textit{Does the association between recalled childhood socioeconomic disadvantage and later-life measured cognitive performance differ across memory, orientation, and the combined memory-orientation score, under defensible scoring rules?}

\section{Methodology}
To resolve denominator inconsistencies found in broad cognitive batteries, we utilized three continuous outcomes administered universally to direct respondents:
\begin{itemize}
    \item \textbf{Memory Score (0--20)}
    \item \textbf{Orientation Score (0--8)}
    \item \textbf{Combined Memory-Orientation Score (0--28)}
\end{itemize}

Linear regression models (Gaussian family) were estimated using survey weights with Taylor-series linearization for clustered standard errors. 
\begin{itemize}
    \item \textbf{Model A (Baseline):} Adjusts for age, sex, caste, religion, residence, and state fixed effects.
    \item \textbf{Model B (Adult SES):} Adds adult education and MPCE quintile.
    \item \textbf{Model C (Exploratory Health):} Adds ADL, self-rated health, and morbidity.
\end{itemize}

\section{Results}
Results were estimated on the common complete-case sample of 57,245 older adults.

\begin{table}[H]
\centering
\small
\caption{\textbf{Continuous Cognitive Sensitivity (Model B: Adult SES Adjusted)}}
\begin{tabular}{lccc}
\toprule
\textbf{Outcome Domain} & \textbf{Poor SES Effect (95\% CI)} & \textbf{\textit{p}-value} & \textbf{Finding} \\
\midrule
Combined Score (0-28) & $-0.319$ ($-0.554, -0.084$) & 0.008 & Persistent Deficit \\
Memory Score (0-20) & $-0.249$ ($-0.449, -0.050$) & 0.014 & Persistent Deficit \\
Orientation Score (0-8) & $-0.069$ ($-0.163, 0.025$) & 0.148 & Null \\
\bottomrule
\end{tabular}
\end{table}

\section{Conclusion}
While the primary report demonstrated that the \textit{binary} risk of cognitive impairment becomes statistically null upon adult-SES adjustment, the continuous score retains a small, persistent penalty. Crucially, domain separation reveals this penalty is driven almost entirely by the memory domain, whereas orientation is fully attenuated by adult socioeconomic factors.

\end{document}
